\subsection{Further restrictions on the actual icosahedral kernel}

\statement{B.50}{Proved}{Every nontrivial Sylow subgroup of the kernel is nonabelian}
Every nontrivial pro-$p$ Sylow subgroup of $\mathcal K$
is a free pro-$p$ group of rank at least two.  Moreover
$\mathcal K$ has no nontrivial closed abelian subgroup
stable under the based $G_{\mathbb Q(\sqrt5)}$ action
of B.31.  In particular $Z(\mathcal K)=1$, and if
$\mathcal K$ has finite derived length then $\mathcal K=1$.
The last assertion concerns solvability with finite
derived length; it does not exclude a nontrivial
prosolvable group of unbounded derived length.

\begin{proof}
Write $k=\mathbb Q(\sqrt5)$ and
$F_T^{\wedge}=\pi_1(C_{T,\overline k},\overline v_0)$.
By B.31, $\mathcal K\rtimes G_k$ is an actual
closed subgroup of the arithmetic fundamental
group $F_T^{\wedge}\rtimes G_k$ for the specified
rational-point section.  If $S$ is a pro-$p$
Sylow subgroup of $\mathcal K$, then every
$g\in G_k$ carries $S$ to another such Sylow
subgroup.  Profinite Sylow conjugacy supplies
$a\in\mathcal K$ with $ag(S)a^{-1}=S$.
Consequently the closed subgroup
\[
       N_{\mathcal K\rtimes G_k}(S)
                        \longrightarrow G_k               \tag{B.50.1}
\]
is surjective.  This is surjectivity of an
actual compact normalizer projection, not
an assertion that the chosen Sylow is
fixed by the original section.

Closed pro-$p$ subgroups of a free profinite
group are free pro-$p$: their continuous
$p$-cohomological dimension is at most one.
If $S\ne1$ had rank one, it would be
$\mathbb Z_p$.  Apply B.29 to $S$ and
the normalizer (B.50.1).  Its normalizer
field is the full totally real field $k$,
so $S$ would lie in a tetrahedral cusp
inertia group.  This contradicts B.32.
Thus every nontrivial $S$ has rank at
least two, and is nonabelian.

For a closed abelian $G_k$-stable subgroup
$A\leq\mathcal K$, the argument of B.33
gives $A\simeq\prod_{p\in I}\mathbb Z_p$.
Its semidirect product with the same
$G_k$ is an arithmetic normalizer.
B.29 and B.32 again give $A=1$.
The center is characteristic, so this
applies to $Z(\mathcal K)$.
If the closed derived series has
finite length and $\mathcal K\ne1$,
its last nontrivial term is a
nontrivial closed abelian characteristic
subgroup, a contradiction.  Alternatively
each Sylow would be a solvable free
pro-$p$ group and hence procyclic:
a free pro-$p$ group of rank at least
two contains a free pro-$p$ group of
rank two, whose dense abstract free
subgroup is nonsolvable.  This gives
the same contradiction to the first
part of the statement.
\end{proof}

\statement{B.51}{Proved}{Complex conjugation excludes the odd part and constrains odd Sylow normalizers}
Let $c\in G_{\mathbb Q(\sqrt5)}$ be complex
conjugation for the chosen real embedding.
Then
\[
 C_{\widehat F_T}(c)=1,\qquad C_{\mathcal K}(c)=1.            \tag{B.51.1}
\]
Every pro-odd closed subgroup of $\mathcal K$
stable under the based $G_{\mathbb Q(\sqrt5)}$
action is trivial.  In particular a nontrivial
$\mathcal K$ has a nonabelian pro-$2$ Sylow
subgroup, and its largest closed normal
pro-odd subgroup $O_{2'}(\mathcal K)$ is
trivial.

For every nontrivial odd-prime Sylow subgroup
$S$ of $\mathcal K$, its geometric normalizer
$N_{\mathcal K}(S)$ has a nontrivial pro-$2$
Sylow subgroup.  If $S$ is finitely generated,
then $C_{\mathcal K}(S)$ contains a nontrivial
pro-$2$ subgroup.  These assertions do not
establish that $S$ has infinite rank.

\begin{proof}
Over $\mathbb R$, rescale the coordinate
of (B.31.1) so that the punctures are
$\infty,i,-i$.  Its real locus is
$\mathbb R$.  Move the real basepoint
to zero along this real locus; that
path is fixed by complex conjugation.
Choose a positive meridian $a$ around
$i$ and its reflected loop $b=c(a)$,
which is a negative meridian around
$-i$.  These two loops form a free
basis and $c$ exchanges them.
The based topological extension is
therefore
\[
 \langle a,b,c\mid c^2=1,\ cac=b,\ cbc=a\rangle
          \simeq\mathbb Z*C_2.                             \tag{B.51.2}
\]
Finite equivariant covering comparison
and the full induced topology on its
index-two free subgroup give
\[
          \widehat F_T\rtimes\langle c\rangle
                     \simeq\widehat{\mathbb Z}\amalg C_2.
                                                               \tag{B.51.3}
\]
This is the real action supplied by
the rational-point section of B.31.
Herfort--Ribes, Theorems A' and B',
\cite{A_HR1985}, give both that the
centralizer of $c$ in (B.51.3) is
$\langle c\rangle$ and that every
involution there is conjugate to
$c$.  Thus (B.51.1) holds, and
indeed every such involution has
trivial centralizer in $\widehat F_T$.

We need the full pro-odd lifting
argument, rather than an inference
from one prime quotient.  Let
$H\leq\mathcal K$ be closed,
pro-odd and $c$-stable.  For every
$c$-stable open normal $V\leq H$,
the fixed-point map
\[
                   C_H(c)\longrightarrow C_{H/V}(c)
\]
is surjective.  Given a fixed coset,
a lift conjugates $\langle c\rangle$
to another complement inside
$V\rtimes\langle c\rangle$.
The two complements are pro-$2$
Sylow subgroups of this group,
since $V$ is pro-odd; profinite
Sylow conjugacy corrects the lift
by an element of $V$.  This is
the claimed fixed-point lift.
Such $c$-stable $V$ are cofinal,
by replacing any open normal
subgroup with its intersection
with its $c$-image.

By (B.51.1), the finite group
$H/V$ has a fixed-point-free
involution.  A finite group
$D$ with such an automorphism
$\alpha$ is abelian and $\alpha$
is inversion.  Indeed
$d\mapsto d^{-1}\alpha(d)$ is
injective, because equality
of two values makes their
ratio fixed; finiteness makes
it surjective.  Every value
is a cocycle, so surjectivity
gives $\alpha(h)=h^{-1}$ for
every $h\in D$.  Since
$\alpha$ is a homomorphism,
inversion must be a homomorphism,
and $D$ is abelian.  All
finite quotients just considered
are therefore abelian and
$c$ inverts them, so $H$ is
abelian and is inverted by $c$.
If $H$ is stable under the
whole based $G_k$, B.50 gives
$H=1$.  Applying this to
$H=\mathcal K$ when it is
pro-odd shows that a nontrivial
$\mathcal K$ has nonzero
$2$-part; B.50 makes its
$2$-Sylow nonabelian.
The subgroup $O_{2'}(\mathcal K)$
is characteristic and hence
$G_k$-stable, giving its
triviality as well.

Now let $p$ be odd and
$1\ne S\in\operatorname{Syl}_p(\mathcal K)$.
Put $D=N_{\mathcal K}(S)$ and
$E_S=N_{\mathcal K\rtimes\langle c\rangle}(S)$.
The same Sylow-conjugacy argument
as (B.50.1) gives an exact sequence
\[
                      1\longrightarrow D\longrightarrow E_S
                        \longrightarrow C_2\longrightarrow1.
                                                               \tag{B.51.4}
\]
There is no involution in $E_S$.
Such an involution would normalize
$S$, and (B.51.3) would make
its fixed subgroup in $S$
trivial.  The preceding coprime
fixed-point argument, now on
the pro-$p$ group $S$, would
make $S$ abelian, contradicting
B.50.  If $D$ were pro-odd,
a $2$-Sylow subgroup of $E_S$
would map isomorphically onto
the quotient $C_2$ of (B.51.4),
producing just such an involution.
Thus $D$ has a nontrivial
pro-$2$ Sylow subgroup $T$.

Finally suppose $S$ has finite
rank $r$.  Its Frattini
automorphism kernel is pro-$p$,
so the pro-$2$ image of $T$
in $\operatorname{Aut}(S)$
injects into the finite group
$\operatorname{GL}_r(\mathbb F_p)$.
The image is finite.  Since
$\mathcal K$ is torsion-free,
$T$ is infinite, and the
kernel of this action is
a nontrivial open subgroup
of $T$.  It centralizes $S$
and proves the last assertion.
\end{proof}

The general free-profinite
normalizer shortcut cannot
upgrade this conclusion to
infinite rank.  For distinct
primes $p,q$ put
\[
               J=\widehat F_2^{(p)}\times\widehat F_2^{(q)}.
                                                               \tag{B.51.5}
\]
It is generated by $(x,u)$
and $(y,v)$: primary powers
of each generator recover
its two components.  Each
nontrivial Sylow subgroup
is free pro-prime of rank
two, so the cohomological
Sylow criterion makes $J$
projective.  The epimorphism
$\widehat F_2\to J$ consequently
has a continuous section,
which embeds $J$ as a closed
retract of $\widehat F_2$.
Its nonabelian $p$-Sylow is
normal and is centralized
by its nonabelian $q$-Sylow.
Thus a finite-rank nonabelian
Sylow in a closed subgroup
of $\widehat F_2$ need not
be self-normalizing.
This is a counterexample
to that general shortcut,
not a model of the actual
arithmetically constrained
$\mathcal K$.  The script
\texttt{check\_normalizer\_model.py}
checks its indicated small
Heisenberg quotient; the
infinite embedding is proved
by projectivity above.

The unresolved upgrade is
to exclude the nontrivial
commuting pro-$2$ subgroup
forced by B.51 when an odd
Sylow has finite rank, using
some additional property of
this specific centralizer.
No general assertion that
all Sylows are $c$-invariant
has been used: (B.51.4)
need not split, and a
pro-odd hypothesis on its
geometric kernel is exactly
what is missing.
